% Original pages: 17--24
% Figures: none
% Uncertainty log: none

\noindent\textit{Proposition 1.11.\amtarget{1.11}} i) \textit{Let $\mathfrak p_1,\ldots,\mathfrak p_n$ be prime ideals and let $\mathfrak a$ be an ideal contained in $\bigcup_{i=1}^n\mathfrak p_i$. Then $\mathfrak a\subseteq\mathfrak p_i$ for some $i$.}

ii) \textit{Let $\mathfrak a_1,\ldots,\mathfrak a_n$ be ideals and let $\mathfrak p$ be a prime ideal containing $\bigcap_{i=1}^n\mathfrak a_i$. Then $\mathfrak p\supseteq\mathfrak a_i$ for some $i$. If $\mathfrak p=\bigcap\mathfrak a_i$, then $\mathfrak p=\mathfrak a_i$ for some $i$.}

\noindent\textit{Proof.} i) is proved by induction on $n$ in the form
\[
 \mathfrak a\nsubseteq\mathfrak p_i\ (1\leq i\leq n)
 \quad\Rightarrow\quad
 \mathfrak a\nsubseteq\bigcup_{i=1}^n\mathfrak p_i.
\]
It is certainly true for $n=1$. If $n>1$ and the result is true for $n-1$, then for each $i$ there exists $x_i\in\mathfrak a$ such that $x_i\notin\mathfrak p_j$ whenever $j\neq i$. If for some $i$ we have $x_i\notin\mathfrak p_i$, we are through. If not, then $x_i\in\mathfrak p_i$ for all $i$. Consider the element
\[
 y=\sum_{i=1}^n x_1x_2\cdots x_{i-1}x_{i+1}x_{i+2}\cdots x_n;
\]
we have $y\in\mathfrak a$ and $y\notin\mathfrak p_i$ $(1\leq i\leq n)$. Hence $\mathfrak a\nsubseteq\bigcup_{i=1}^n\mathfrak p_i$.

ii) Suppose $\mathfrak p\nsupseteq\mathfrak a_i$ for all $i$. Then there exist $x_i\in\mathfrak a_i$, $x_i\notin\mathfrak p$ $(1\leq i\leq n)$, and therefore $\prod x_i\in\prod\mathfrak a_i\subseteq\bigcap\mathfrak a_i$; but $\prod x_i\notin\mathfrak p$ (since $\mathfrak p$ is prime). Hence $\mathfrak p\nsupseteq\bigcap\mathfrak a_i$. Finally, if $\mathfrak p=\bigcap\mathfrak a_i$, then $\mathfrak p\subseteq\mathfrak a_i$ and hence $\mathfrak p=\mathfrak a_i$ for some $i$. \proofbox

If $\mathfrak a$, $\mathfrak b$ are ideals in a ring $A$, their \textit{ideal quotient} is
\[
 (\mathfrak a:\mathfrak b)=\{x\in A:x\mathfrak b\subseteq\mathfrak a\}
\]
which is an ideal. In particular, $(0:\mathfrak b)$ is called the \textit{annihilator} of $\mathfrak b$ and is also denoted by $\operatorname{Ann}(\mathfrak b)$: it is the set of all $x\in A$ such that $x\mathfrak b=0$. In this notation the set of all zero-divisors in $A$ is
\[
 D=\bigcup_{x\neq0}\operatorname{Ann}(x).
\]
If $\mathfrak b$ is a principal ideal $(x)$, we shall write $(\mathfrak a:x)$ in place of $(\mathfrak a:(x))$.

\textline
\noindent\textbf{Examples.} 1) In $\mathbb Z$,
\[
((12):(8))=(3),
\]
because $12$ divides $8a$ precisely when $3$ divides $a$.  More generally,
$((m):(n))=(m/(m,n))$.

2) In $k[x]$ we have $((x^3):(x))=(x^2)$.  In $k[x,y]$ we have
\[
((x^2,xy):(x))=(x,y),
\]
because $fx$ lies in $(x^2,xy)=x(x,y)$ precisely when $f$ lies in $(x,y)$.

3) In $\mathbb Z/12\mathbb Z$, the annihilator of the class of $4$ is the
ideal generated by the class of $3$.  In
$k[\varepsilon]/(\varepsilon^2)$, the annihilator of $\varepsilon$ is
$(\varepsilon)$.  In an integral domain, the annihilator of every non-zero
element is zero.
\textline

\noindent\textbf{Example.} If $A=\mathbb Z$, $\mathfrak a=(m)$, $\mathfrak b=(n)$, where say $m=\prod_p p^{\mu_p}$, $n=\prod_p p^{\nu_p}$, then $(\mathfrak a:\mathfrak b)=(q)$ where $q=\prod_p p^{\gamma_p}$ and
\[
 \gamma_p=\max(\mu_p-\nu_p,0)=\mu_p-\min(\mu_p,\nu_p).
\]
Hence $q=m/(m,n)$, where $(m,n)$ is the h.c.f. of $m$ and $n$.

\noindent\textit{Exercise 1.12.\amtarget{1.12}}
\begin{enumerate}
\renewcommand{\labelenumi}{\roman{enumi})}
\item $\mathfrak a\subseteq(\mathfrak a:\mathfrak b)$
\item $(\mathfrak a:\mathfrak b)\mathfrak b\subseteq\mathfrak a$
\item $((\mathfrak a:\mathfrak b):\mathfrak c)=(\mathfrak a:\mathfrak b\mathfrak c)=((\mathfrak a:\mathfrak c):\mathfrak b)$
\item $(\bigcap_i\mathfrak a_i:\mathfrak b)=\bigcap_i(\mathfrak a_i:\mathfrak b)$
\item $(\mathfrak a:\sum_i\mathfrak b_i)=\bigcap_i(\mathfrak a:\mathfrak b_i)$.
\end{enumerate}

If $\mathfrak a$ is any ideal of $A$, the \textit{radical} of $\mathfrak a$ is
\[
 r(\mathfrak a)=\{x\in A:x^n\in\mathfrak a\text{ for some }n>0\}.
\]
If $\phi:A\longrightarrow A/\mathfrak a$ is the standard homomorphism, then $r(\mathfrak a)=\phi^{-1}(\mathfrak N_{A/\mathfrak a})$ and hence $r(\mathfrak a)$ is an ideal by \amref{1.7}.

\textline
\noindent\textbf{Examples.} 1) In $\mathbb Z$,
\[
r((12))=(6),
\]
since an integer has a power divisible by $12$ precisely when it is divisible
by both $2$ and $3$.  More generally, the radical of $(n)$ is generated by the
product of the distinct prime divisors of $n$.

2) In $k[x]$, an element has a power in $(x^4)$ precisely when it is divisible
by $x$, so
\[
r((x^4))=(x).
\]
Likewise, $r((f^n))=(f)$ when $f$ is irreducible and $n>0$.

3) In $k[x,y]$,
\[
r((x^2,y^3))=(x,y).
\]
Indeed, $x$ and $y$ belong to the radical, so $(x,y)$ is contained in it; the
reverse inclusion follows because $(x,y)$ is maximal and contains
$(x^2,y^3)$.
\textline

\noindent\textit{Exercise 1.13.\amtarget{1.13}}
\begin{enumerate}
\renewcommand{\labelenumi}{\roman{enumi})}
\item $r(\mathfrak a)\supseteq\mathfrak a$
\item $r(r(\mathfrak a))=r(\mathfrak a)$
\item $r(\mathfrak a\mathfrak b)=r(\mathfrak a\cap\mathfrak b)=r(\mathfrak a)\cap r(\mathfrak b)$
\item $r(\mathfrak a)=(1)\Longleftrightarrow\mathfrak a=(1)$
\item $r(\mathfrak a+\mathfrak b)=r(r(\mathfrak a)+r(\mathfrak b))$
\item if $\mathfrak p$ is prime, $r(\mathfrak p^n)=\mathfrak p$ for all $n>0$.
\end{enumerate}

\noindent\textit{Proposition 1.14.\amtarget{1.14} The radical of an ideal $\mathfrak a$ is the intersection of the prime ideals which contain $\mathfrak a$.}

\noindent\textit{Proof.} Apply \amref{1.8} to $A/\mathfrak a$. \proofbox

More generally, we may define the radical $r(E)$ of any subset $E$ of $A$ in the same way. It is \textit{not} an ideal in general. We have $r(\bigcup_\alpha E_\alpha)=\bigcup_\alpha r(E_\alpha)$, for any family of subsets $E_\alpha$ of $A$.

\noindent\textit{Proposition 1.15.\amtarget{1.15} $D={}$ set of zero-divisors of $A=\bigcup_{x\neq0}r(\operatorname{Ann}(x))$.}

\noindent\textit{Proof.} $D=r(D)=r(\bigcup_{x\neq0}\operatorname{Ann}(x))=\bigcup_{x\neq0}r(\operatorname{Ann}(x))$. \proofbox

\noindent\textbf{Example.} If $A=\mathbb Z$, $\mathfrak a=(m)$, let $p_i$ $(1\leq i\leq r)$ be the distinct prime divisors of $m$. Then $r(\mathfrak a)=(p_1\cdots p_r)=\bigcap_{i=1}^r(p_i)$.

\noindent\textit{Proposition 1.16.\amtarget{1.16} Let $\mathfrak a$, $\mathfrak b$ be ideals in a ring $A$ such that $r(\mathfrak a)$, $r(\mathfrak b)$ are coprime. Then $\mathfrak a$, $\mathfrak b$ are coprime.}

\noindent\textit{Proof.} $r(\mathfrak a+\mathfrak b)=r(r(\mathfrak a)+r(\mathfrak b))=r(1)=(1)$, hence $\mathfrak a+\mathfrak b=(1)$ by \amref{1.13}.

\subsection{Extension and contraction}

Let $f:A\longrightarrow B$ be a ring homomorphism. If $\mathfrak a$ is an ideal in $A$, the set $f(\mathfrak a)$ is not necessarily an ideal in $B$ (e.g., let $f$ be the embedding of $\mathbb Z$ in $\mathbb Q$, the field of rationals, and take $\mathfrak a$ to be any non-zero ideal in $\mathbb Z$.) We define the \textit{extension} $\mathfrak a^e$ of $\mathfrak a$ to be the ideal $Bf(\mathfrak a)$ generated by $f(\mathfrak a)$ in $B$: explicitly, $\mathfrak a^e$ is the set of all sums $\sum y_i f(x_i)$ where $x_i\in\mathfrak a$, $y_i\in B$.

If $\mathfrak b$ is an ideal of $B$, then $f^{-1}(\mathfrak b)$ is always an ideal of $A$, called the \textit{contraction} $\mathfrak b^c$ of $\mathfrak b$. If $\mathfrak b$ is prime, then $\mathfrak b^c$ is prime. If $\mathfrak a$ is prime, $\mathfrak a^e$ need not be prime (for example, $f:\mathbb Z\longrightarrow\mathbb Q$, $\mathfrak a\neq0$; then $\mathfrak a^e=\mathbb Q$, which is not a prime ideal).

\textline
\noindent\textbf{Examples.} 1) For the inclusion
$k[x]\longrightarrow k[x,y]$, the extensions of $(x)$ and $(x^2)$ are the
ideals $(x)$ and $(x^2)$ of $k[x,y]$.  The contractions of $(x,y)$ and
$(x^2,y)$ are $(x)$ and $(x^2)$, respectively.

2) For the quotient map $\pi:\mathbb Z\longrightarrow\mathbb Z/6\mathbb Z$,
the extension of $(2)$ is the ideal generated by the class of $2$, and its
contraction is again $(2)$.  The extension of $(5)$ is the unit ideal because
the class of $5$ is a unit; its contraction is therefore all of $\mathbb Z$.

3) For the inclusion $\mathbb Z\longrightarrow\mathbb Q$, every non-zero ideal
of $\mathbb Z$ extends to the unit ideal of $\mathbb Q$, while $(0)$ extends to
$(0)$.  The only ideals of the field $\mathbb Q$ are $(0)$ and $(1)$, and their
contractions are $(0)$ and $\mathbb Z$.
\textline

We can factorize $f$ as follows:
\[
 A\xrightarrow{p}f(A)\xrightarrow{j}B
\]
where $p$ is surjective and $j$ is injective. For $p$ the situation is very simple \amref{1.1}: there is a one-to-one correspondence between ideals of $f(A)$ and ideals of $A$ which contain $\operatorname{Ker}(f)$, and prime ideals correspond to prime ideals. For $j$, on the other hand, the general situation is very complicated. The classical example is from algebraic number theory.

\noindent\textbf{Example.} Consider $\mathbb Z\longrightarrow\mathbb Z[i]$, where $i=\sqrt{-1}$. A prime ideal $(p)$ of $\mathbb Z$ may or may not stay prime when extended to $\mathbb Z[i]$. In fact $\mathbb Z[i]$ is a principal ideal domain (because it has a Euclidean algorithm) and the situation is as follows:
\begin{enumerate}
\renewcommand{\labelenumi}{\roman{enumi})}
\item $(2)^e=((1+i)^2)$, the square of a prime ideal in $\mathbb Z[i]$;
\item If $p\equiv1\pmod 4$ then $(p)^e$ is the product of two distinct prime ideals (for example, $(5)^e=(2+i)(2-i)$);
\item If $p\equiv3\pmod 4$ then $(p)^e$ is prime in $\mathbb Z[i]$.
\end{enumerate}
Of these, ii) is not a trivial result. It is effectively equivalent to a theorem of Fermat which says that a prime $p\equiv1\pmod 4$ can be expressed, essentially uniquely, as a sum of two integer squares (thus $5=2^2+1^2$, $97=9^2+4^2$, etc.).

In fact the behavior of prime ideals under extensions of this sort is one of the central problems of algebraic number theory.

Let $f:A\longrightarrow B$, $\mathfrak a$ and $\mathfrak b$ be as before. Then

\noindent\textit{Proposition 1.17.\amtarget{1.17}}
\begin{enumerate}
\renewcommand{\labelenumi}{\roman{enumi})}
\item $\mathfrak a\subseteq\mathfrak a^{ec}$, $\mathfrak b\supseteq\mathfrak b^{ce}$;
\item $\mathfrak b^c=\mathfrak b^{cec}$, $\mathfrak a^e=\mathfrak a^{ece}$;
\item \textit{If $C$ is the set of contracted ideals in $A$ and if $E$ is the set of extended ideals in $B$, then $C=\{\mathfrak a\mid\mathfrak a^{ec}=\mathfrak a\}$, $E=\{\mathfrak b\mid\mathfrak b^{ce}=\mathfrak b\}$, and $\mathfrak a\mapsto\mathfrak a^e$ is a bijective map of $C$ onto $E$, whose inverse is $\mathfrak b\mapsto\mathfrak b^c$.}
\end{enumerate}

\noindent\textit{Proof.} i) is trivial, and ii) follows from i).

iii) If $\mathfrak a\in C$, then $\mathfrak a=\mathfrak b^c=\mathfrak b^{cec}=\mathfrak a^{ec}$; conversely if $\mathfrak a=\mathfrak a^{ec}$ then $\mathfrak a$ is the contraction of $\mathfrak a^e$. Similarly for $E$. \proofbox

\noindent\textit{Exercise 1.18.\amtarget{1.18} If $\mathfrak a_1$, $\mathfrak a_2$ are ideals of $A$ and if $\mathfrak b_1$, $\mathfrak b_2$ are ideals of $B$, then}
\[
\begin{array}{ll}
(\mathfrak a_1+\mathfrak a_2)^e=\mathfrak a_1^e+\mathfrak a_2^e,
& (\mathfrak b_1+\mathfrak b_2)^c\supseteq\mathfrak b_1^c+\mathfrak b_2^c,\\
(\mathfrak a_1\cap\mathfrak a_2)^e\subseteq\mathfrak a_1^e\cap\mathfrak a_2^e,
& (\mathfrak b_1\cap\mathfrak b_2)^c=\mathfrak b_1^c\cap\mathfrak b_2^c,\\
(\mathfrak a_1\mathfrak a_2)^e=\mathfrak a_1^e\mathfrak a_2^e,
& (\mathfrak b_1\mathfrak b_2)^c\supseteq\mathfrak b_1^c\mathfrak b_2^c,\\
(\mathfrak a_1:\mathfrak a_2)^e\subseteq(\mathfrak a_1^e:\mathfrak a_2^e),
& (\mathfrak b_1:\mathfrak b_2)^c\subseteq(\mathfrak b_1^c:\mathfrak b_2^c),\\
r(\mathfrak a)^e\subseteq r(\mathfrak a^e),
& r(\mathfrak b)^c=r(\mathfrak b^c).
\end{array}
\]
\textit{The set of ideals $E$ is closed under sum and product, and $C$ is closed under the other three operations.}

\subsection{Exercises}

\begin{enumerate}
\amstaritem Let $x$ be a nilpotent element of a ring $A$. Show that $1+x$ is a unit of $A$. Deduce that the sum of a nilpotent element and a unit is a unit.

\amstaritem Let $A$ be a ring and let $A[x]$ be the ring of polynomials in an indeterminate $x$, with coefficients in $A$. Let $f=a_0+a_1x+\cdots+a_nx^n\in A[x]$. Prove that
\begin{enumerate}
\renewcommand{\labelenumii}{\roman{enumii})}
\item $f$ is a unit in $A[x]\Longleftrightarrow a_0$ is a unit in $A$ and $a_1,\ldots,a_n$ are nilpotent. [If $b_0+b_1x+\cdots+b_mx^m$ is the inverse of $f$, prove by induction on $r$ that $a_n^{r+1}b_{m-r}=0$. Hence show that $a_n$ is nilpotent, and then use Ex. 1.]
\item $f$ is nilpotent $\Longleftrightarrow a_0,a_1,\ldots,a_n$ are nilpotent.
\item $f$ is a zero-divisor $\Longleftrightarrow$ there exists $a\neq0$ in $A$ such that $af=0$. [Choose a polynomial $g=b_0+b_1x+\cdots+b_mx^m$ of least degree $m$ such that $fg=0$. Then $a_nb_m=0$, hence $a_ng=0$ (because $a_ng$ annihilates $f$ and has degree $<m$). Now show by induction that $a_{n-r}g=0$ $(0\leq r\leq n)$.]
\item $f$ is said to be \textit{primitive} if $(a_0,a_1,\ldots,a_n)=(1)$. Prove that if $f,g\in A[x]$, then $fg$ is primitive $\Longleftrightarrow f$ and $g$ are primitive.
\end{enumerate}

\item Generalize the results of Exercise 2 to a polynomial ring $A[x_1,\ldots,x_r]$ in several indeterminates.

\amstaritem In the ring $A[x]$, the Jacobson radical is equal to the nilradical.

\item Let $A$ be a ring and let $A[[x]]$ be the ring of \textit{formal power series} $f=\sum_{n=0}^\infty a_nx^n$ with coefficients in $A$. Show that
\begin{enumerate}
\renewcommand{\labelenumii}{\roman{enumii})}
\item $f$ is a unit in $A[[x]]\Longleftrightarrow a_0$ is a unit in $A$.
\item If $f$ is nilpotent, then $a_n$ is nilpotent for all $n\geq0$. Is the converse true? (See Chapter 7, Exercise 2.)
\item $f$ belongs to the Jacobson radical of $A[[x]]\Longleftrightarrow a_0$ belongs to the Jacobson radical of $A$.
\item The contraction of a maximal ideal $\mathfrak m$ of $A[[x]]$ is a maximal ideal of $A$, and $\mathfrak m$ is generated by $\mathfrak m^c$ and $x$.
\item Every prime ideal of $A$ is the contraction of a prime ideal of $A[[x]]$.
\end{enumerate}

\amstaritem A ring $A$ is such that every ideal not contained in the nilradical contains a non-zero idempotent (that is, an element $e$ such that $e^2=e\neq0$). Prove that the nilradical and Jacobson radical of $A$ are equal.

\item Let $A$ be a ring in which every element $x$ satisfies $x^n=x$ for some $n>1$ (depending on $x$). Show that every prime ideal in $A$ is maximal.

\amstaritem Let $A$ be a ring $\neq0$. Show that the set of prime ideals of $A$ has minimal elements with respect to inclusion.

\amstaritem Let $\mathfrak a$ be an ideal $\neq(1)$ in a ring $A$. Show that $\mathfrak a=r(\mathfrak a)\Longleftrightarrow\mathfrak a$ is an intersection of prime ideals.

\item Let $A$ be a ring, $\mathfrak N$ its nilradical. Show that the following are equivalent:
\begin{enumerate}
\renewcommand{\labelenumii}{\roman{enumii})}
\item $A$ has exactly one prime ideal;
\item every element of $A$ is either a unit or nilpotent;
\item $A/\mathfrak N$ is a field.
\end{enumerate}

\item A ring $A$ is \textit{Boolean} if $x^2=x$ for all $x\in A$. In a Boolean ring $A$, show that
\begin{enumerate}
\renewcommand{\labelenumii}{\roman{enumii})}
\item $2x=0$ for all $x\in A$;
\item every prime ideal $\mathfrak p$ is maximal, and $A/\mathfrak p$ is a field with two elements;
\item every finitely generated ideal in $A$ is principal.
\end{enumerate}

\item A local ring contains no idempotent $\neq0,1$.

\noindent\textit{Construction of an algebraic closure of a field} (E. Artin).

\item Let $K$ be a field and let $\Sigma$ be the set of all irreducible monic polynomials $f$ in one indeterminate with coefficients in $K$. Let $A$ be the polynomial ring over $K$ generated by indeterminates $x_f$, one for each $f\in\Sigma$. Let $\mathfrak a$ be the ideal of $A$ generated by the polynomials $f(x_f)$ for all $f\in\Sigma$. Show that $\mathfrak a\neq(1)$.

Let $\mathfrak m$ be a maximal ideal of $A$ containing $\mathfrak a$, and let $K_1=A/\mathfrak m$. Then $K_1$ is an extension field of $K$ in which each $f\in\Sigma$ has a root. Repeat the construction with $K_1$ in place of $K$, obtaining a field $K_2$, and so on. Let $L=\bigcup_{n=1}^\infty K_n$. Then $L$ is a field in which each $f\in\Sigma$ splits completely into linear factors. Let $\overline K$ be the set of all elements of $L$ which are algebraic over $K$. Then $\overline K$ is an algebraic closure of $K$.

\item In a ring $A$, let $\Sigma$ be the set of all ideals in which every element is a zero-divisor. Show that the set $\Sigma$ has maximal elements and that every maximal element of $\Sigma$ is a prime ideal. Hence the set of zero-divisors in $A$ is a union of prime ideals.

\noindent\textit{The prime spectrum of a ring}

\item Let $A$ be a ring and let $X$ be the set of all prime ideals of $A$. For each subset $E$ of $A$, let $V(E)$ denote the set of all prime ideals of $A$ which contain $E$. Prove that
\begin{enumerate}
\renewcommand{\labelenumii}{\roman{enumii})}
\item if $\mathfrak a$ is the ideal generated by $E$, then $V(E)=V(\mathfrak a)=V(r(\mathfrak a))$.
\item $V(0)=X$, $V(1)=\varnothing$.
\item if $(E_i)_{i\in I}$ is any family of subsets of $A$, then
\[
 V\left(\bigcup_{i\in I}E_i\right)=\bigcap_{i\in I}V(E_i).
\]
\item $V(\mathfrak a\cap\mathfrak b)=V(\mathfrak a\mathfrak b)=V(\mathfrak a)\cup V(\mathfrak b)$ for any ideals $\mathfrak a$, $\mathfrak b$ of $A$.
\end{enumerate}
These results show that the sets $V(E)$ satisfy the axioms for closed sets in a topological space. The resulting topology is called the \textit{Zariski topology}. The topological space $X$ is called the \textit{prime spectrum} of $A$, and is written $\operatorname{Spec}(A)$.

\item Draw pictures of $\operatorname{Spec}(\mathbb Z)$, $\operatorname{Spec}(\mathbb R)$, $\operatorname{Spec}(\mathbb C[x])$, $\operatorname{Spec}(\mathbb R[x])$, $\operatorname{Spec}(\mathbb Z[x])$.

\item For each $f\in A$, let $X_f$ denote the complement of $V(f)$ in $X=\operatorname{Spec}(A)$. The sets $X_f$ are open. Show that they form a basis of open sets for the Zariski topology, and that
\begin{enumerate}
\renewcommand{\labelenumii}{\roman{enumii})}
\item $X_f\cap X_g=X_{fg}$;
\item $X_f=\varnothing\Longleftrightarrow f$ is nilpotent;
\item $X_f=X\Longleftrightarrow f$ is a unit;
\item $X_f=X_g\Longleftrightarrow r((f))=r((g))$;
\item $X$ is quasi-compact (that is, every open covering of $X$ has a finite sub-covering).
\item More generally, each $X_f$ is quasi-compact.
\item An open subset of $X$ is quasi-compact if and only if it is a finite union of sets $X_f$.
\end{enumerate}
The sets $X_f$ are called \textit{basic open sets} of $X=\operatorname{Spec}(A)$.

[To prove (v), remark that it is enough to consider a covering of $X$ by basic open sets $X_{f_i}$ $(i\in I)$. Show that the $f_i$ generate the unit ideal and hence that there is an equation of the form
\[
 1=\sum_{i\in J}g_if_i\qquad(g_i\in A)
\]
where $J$ is some finite subset of $I$. Then the $X_{f_i}$ $(i\in J)$ cover $X$.]

\item For psychological reasons it is sometimes convenient to denote a prime ideal of $A$ by a letter such as $x$ or $y$ when thinking of it as a point of $X=\operatorname{Spec}(A)$. When thinking of $x$ as a prime ideal of $A$, we denote it by $\mathfrak p_x$ (logically, of course, it is the same thing). Show that
\begin{enumerate}
\renewcommand{\labelenumii}{\roman{enumii})}
\item the set $\{x\}$ is closed (we say that $x$ is a ``closed point'') in $\operatorname{Spec}(A)\Longleftrightarrow\mathfrak p_x$ is maximal;
\item $\overline{\{x\}}=V(\mathfrak p_x)$;
\item $y\in\overline{\{x\}}\Longleftrightarrow\mathfrak p_x\subseteq\mathfrak p_y$;
\item $X$ is a $T_0$-space (this means that if $x$, $y$ are distinct points of $X$, then either there is a neighborhood of $x$ which does not contain $y$, or else there is a neighborhood of $y$ which does not contain $x$).
\end{enumerate}

\item A topological space $X$ is said to be \textit{irreducible} if $X\neq\varnothing$ and if every pair of non-empty open sets in $X$ intersect, or equivalently if every non-empty open set is dense in $X$. Show that $\operatorname{Spec}(A)$ is irreducible if and only if the nilradical of $A$ is a prime ideal.

\item Let $X$ be a topological space.
\begin{enumerate}
\renewcommand{\labelenumii}{\roman{enumii})}
\item If $Y$ is an irreducible (Exercise 19) subspace of $X$, then the closure $\overline Y$ of $Y$ in $X$ is irreducible.
\item Every irreducible subspace of $X$ is contained in a maximal irreducible subspace.
\item The maximal irreducible subspaces of $X$ are closed and cover $X$. They are called the \textit{irreducible components} of $X$. What are the irreducible components of a Hausdorff space?
\item If $A$ is a ring and $X=\operatorname{Spec}(A)$, then the irreducible components of $X$ are the closed sets $V(\mathfrak p)$, where $\mathfrak p$ is a minimal prime ideal of $A$ (Exercise 8).
\end{enumerate}

\item Let $\phi:A\longrightarrow B$ be a ring homomorphism. Let $X=\operatorname{Spec}(A)$ and $Y=\operatorname{Spec}(B)$. If $\mathfrak q\in Y$, then $\phi^{-1}(\mathfrak q)$ is a prime ideal of $A$, i.e., a point of $X$. Hence $\phi$ induces a mapping $\phi^*:Y\longrightarrow X$. Show that
\begin{enumerate}
\renewcommand{\labelenumii}{\roman{enumii})}
\item If $f\in A$ then $\phi^{*-1}(X_f)=Y_{\phi(f)}$, and hence that $\phi^*$ is continuous.
\item If $\mathfrak a$ is an ideal of $A$, then $\phi^{*-1}(V(\mathfrak a))=V(\mathfrak a^e)$.
\item If $\mathfrak b$ is an ideal of $B$, then $\overline{\phi^*(V(\mathfrak b))}=V(\mathfrak b^c)$.
\item If $\phi$ is surjective, then $\phi^*$ is a homeomorphism of $Y$ onto the closed subset $V(\operatorname{Ker}(\phi))$ of $X$. (In particular, $\operatorname{Spec}(A)$ and $\operatorname{Spec}(A/\mathfrak N)$ (where $\mathfrak N$ is the nilradical of $A$) are naturally homeomorphic.)
\item If $\phi$ is injective, then $\phi^*(Y)$ is dense in $X$. More precisely, $\phi^*(Y)$ is dense in $X\Longleftrightarrow\operatorname{Ker}(\phi)\subseteq\mathfrak N$.
\item Let $\psi:B\longrightarrow C$ be another ring homomorphism. Then $(\psi\circ\phi)^*=\phi^*\circ\psi^*$.
\item Let $A$ be an integral domain with just one non-zero prime ideal $\mathfrak p$, and let $K$ be the field of fractions of $A$. Let $B=(A/\mathfrak p)\times K$. Define $\phi:A\longrightarrow B$ by $\phi(x)=(\bar x,x)$, where $\bar x$ is the image of $x$ in $A/\mathfrak p$. Show that $\phi^*$ is bijective but not a homeomorphism.
\end{enumerate}

\item Let $A=\prod_{i=1}^n A_i$ be the direct product of rings $A_i$. Show that $\operatorname{Spec}(A)$ is the disjoint union of open (and closed) subspaces $X_i$, where $X_i$ is canonically homeomorphic with $\operatorname{Spec}(A_i)$.

Conversely, let $A$ be any ring. Show that the following statements are equivalent:
\begin{enumerate}
\renewcommand{\labelenumii}{\roman{enumii})}
\item $X=\operatorname{Spec}(A)$ is disconnected.
\item $A\cong A_1\times A_2$ where neither of the rings $A_1$, $A_2$ is the zero ring.
\item $A$ contains an idempotent $\neq0,1$.
\end{enumerate}
In particular, the spectrum of a local ring is always connected (Exercise 12).

\item Let $A$ be a Boolean ring (Exercise 11), and let $X=\operatorname{Spec}(A)$.
\begin{enumerate}
\renewcommand{\labelenumii}{\roman{enumii})}
\item For each $f\in A$, the set $X_f$ (Exercise 17) is both open and closed in $X$.
\item Let $f_1,\ldots,f_n\in A$. Show that $X_{f_1}\cup\cdots\cup X_{f_n}=X_f$ for some $f\in A$.
\item The sets $X_f$ are the only subsets of $X$ which are both open and closed. [Let $Y\subseteq X$ be both open and closed. Since $Y$ is open, it is a union of basic open sets $X_f$. Since $Y$ is closed and $X$ is quasi-compact (Exercise 17), $Y$ is quasi-compact. Hence $Y$ is a finite union of basic open sets; now use (ii) above.]
\item $X$ is a compact Hausdorff space.
\end{enumerate}

\item Let $L$ be a lattice, in which the sup and inf of two elements $a$, $b$ are denoted by $a\vee b$ and $a\wedge b$ respectively. $L$ is a \textit{Boolean lattice} (or \textit{Boolean algebra}) if
\begin{enumerate}
\renewcommand{\labelenumii}{\roman{enumii})}
\item $L$ has a least element and a greatest element (denoted by $0$, $1$ respectively).
\item Each of $\vee$, $\wedge$ is distributive over the other.
\item Each $a\in L$ has a unique ``complement'' $a'\in L$ such that $a\vee a'=1$ and $a\wedge a'=0$.
\end{enumerate}
(For example, the set of all subsets of a set, ordered by inclusion, is a Boolean lattice.)

Let $L$ be a Boolean lattice. Define addition and multiplication in $L$ by the rules
\[
 a+b=(a\wedge b')\vee(a'\wedge b),\qquad ab=a\wedge b.
\]
Verify that in this way $L$ becomes a Boolean ring, say $A(L)$.

Conversely, starting from a Boolean ring $A$, define an ordering on $A$ as follows: $a\leq b$ means that $a=ab$. Show that, with respect to this ordering, $A$ is a Boolean lattice. [The sup and inf are given by $a\vee b=a+b+ab$ and $a\wedge b=ab$, and the complement by $a'=1-a$.] In this way we obtain a one-to-one correspondence between (isomorphism classes of) Boolean rings and (isomorphism classes of) Boolean lattices.

\item From the last two exercises deduce Stone's theorem, that every Boolean lattice is isomorphic to the lattice of open-and-closed subsets of some compact Hausdorff topological space.

\item Let $A$ be a ring. The subspace of $\operatorname{Spec}(A)$ consisting of the \textit{maximal ideals} of $A$, with the induced topology, is called the \textit{maximal spectrum} of $A$ and is denoted by $\operatorname{Max}(A)$. For arbitrary commutative rings it does not have the nice functorial properties of $\operatorname{Spec}(A)$ (see Exercise 21), because the inverse image of a maximal ideal under a ring homomorphism need not be maximal.

Let $X$ be a compact Hausdorff space and let $C(X)$ denote the ring of all real-valued continuous functions on $X$ (add and multiply functions by adding and multiplying their values). For each $x\in X$, let $\mathfrak m_x$ be the set of all $f\in C(X)$ such that $f(x)=0$. The ideal $\mathfrak m_x$ is maximal, because it is the kernel of the (surjective) homomorphism $C(X)\longrightarrow\mathbb R$ which takes $f$ to $f(x)$. If $\widetilde X$ denotes $\operatorname{Max}(C(X))$, we have therefore defined a mapping $\mu:X\longrightarrow\widetilde X$, namely $x\mapsto\mathfrak m_x$.

We shall show that $\mu$ is a homeomorphism of $X$ onto $\widetilde X$.
\begin{enumerate}
\renewcommand{\labelenumii}{\roman{enumii})}
\item Let $\mathfrak m$ be any maximal ideal of $C(X)$, and let $V=V(\mathfrak m)$ be the set of common zeros of the functions in $\mathfrak m$: that is,
\[
 V=\{x\in X:f(x)=0\text{ for all }f\in\mathfrak m\}.
\]
Suppose that $V$ is empty. Then for each $x\in X$ there exists $f_x\in\mathfrak m$ such that $f_x(x)\neq0$. Since $f_x$ is continuous, there is an open neighborhood $U_x$ of $x$ in $X$ on which $f_x$ does not vanish. By compactness a finite number of the neighborhoods, say $U_{x_1},\ldots,U_{x_n}$, cover $X$. Let
\[
 f=f_{x_1}^2+\cdots+f_{x_n}^2.
\]
Then $f$ does not vanish at any point of $X$, hence is a unit in $C(X)$. But this contradicts $f\in\mathfrak m$, hence $V$ is not empty.

Let $x$ be a point of $V$. Then $\mathfrak m\subseteq\mathfrak m_x$, hence $\mathfrak m=\mathfrak m_x$ because $\mathfrak m$ is maximal. Hence $\mu$ is surjective.

\item By Urysohn's lemma (this is the only non-trivial fact required in the argument) the continuous functions separate the points of $X$. Hence $x\neq y\Rightarrow\mathfrak m_x\neq\mathfrak m_y$, and therefore $\mu$ is injective.

\item Let $f\in C(X)$; let
\[
 U_f=\{x\in X:f(x)\neq0\}
\]
and let
\[
 \widetilde U_f=\{\mathfrak m\in\widetilde X:f\notin\mathfrak m\}.
\]
Show that $\mu(U_f)=\widetilde U_f$. The open sets $U_f$ (resp. $\widetilde U_f$) form a basis of the topology of $X$ (resp. $\widetilde X$) and therefore $\mu$ is a homeomorphism.

Thus $X$ can be reconstructed from the ring of functions $C(X)$.
\end{enumerate}

\noindent\textit{Affine algebraic varieties}

\item Let $k$ be an algebraically closed field and let
\[
 f_\alpha(t_1,\ldots,t_n)=0
\]
be a set of polynomial equations in $n$ variables with coefficients in $k$. The set $X$ of all points $x=(x_1,\ldots,x_n)\in k^n$ which satisfy these equations is an \textit{affine algebraic variety}.

Consider the set of all polynomials $g\in k[t_1,\ldots,t_n]$ with the property that $g(x)=0$ for all $x\in X$. This set is an ideal $I(X)$ in the polynomial ring, and is called the \textit{ideal of the variety} $X$. The quotient ring
\[
 P(X)=k[t_1,\ldots,t_n]/I(X)
\]
is the ring of polynomial functions on $X$, because two polynomials $g$, $h$ define the same polynomial function on $X$ if and only if $g-h$ vanishes at every point of $X$, that is, if and only if $g-h\in I(X)$.
\end{enumerate}
